% introduction.tex — Paper 8: Arrow of Time and Complexification
% ASSERT_CONVENTION: natural_units=kB_equals_1, entropy_base=nats, jordan_product=(1/2)(ab+ba), clifford_convention=euclidean_positive, complex_structure=u_equals_e7, spin_representation=S10plus_boyle, pati_salam=SU4xSU2LxSU2R

%----------------------------------------------------------------------
\section{Introduction}\label{sec:intro}
%----------------------------------------------------------------------

Why does particle physics use complex numbers?  The Standard Model is
a quantum field theory over $\mathbb{C}$: its gauge group $\SMgauge$
acts on complex representations, and its chiral fermions are
complex Weyl spinors.  The exceptional Jordan algebra
$\hthree$---the $3\times 3$ Hermitian matrices over the octonions
$\mathbb{O}$---provides a natural algebraic setting in which the
Standard Model gauge group and its chiral representation
emerge~\cite{JvNW1934,TodorovDrenska2018,Furey2018,Boyle2020,BGW2020}.
In a companion paper~\cite{Paper7}, we showed that a single algebraic
choice---a unit imaginary octonion $u \in S^6 \subset
\mathrm{Im}(\mathbb{O})$---simultaneously determines both the gauge
group $\SMgauge$ (via the $\Ffour$ intersection of
Todorov--Drenska~\cite{TodorovDrenska2018}) and the chiral
(left-handed) fermion representation (via the $\Cl{6}$ volume form of
Furey~\cite{Furey2018}).  This ``one choice, two consequences'' result
connects the self-modeling axioms of Paper~5~\cite{Paper5} through
$\hthree$ to a one-generation chiral Standard Model representation in
a nine-link chain with explicit gap identification.

The derivation, however, is conditional on a step that Paper~7 left
open.  The Peirce decomposition of $\hthree$ under a rank-1
idempotent $E_{11}$ produces a half-space $\Vhalf \cong \mathbb{O}^2$
that is a \emph{real} 16-dimensional space carrying the spinor
representation of $\Spin{9}$.  The gauge group and chirality
results require this space to be complexified:
$\Vhalf \otimes_{\mathbb{R}} \mathbb{C} = S_{10}^+$, upgrading
$\Spin{9} \to \Spin{10}$ and $\Ffour \to \Esix$.  Paper~7
argued that the C*-algebra nature of the observer (from
Paper~5~\cite{Paper5}) motivates this complexification, but flagged
it as an open gap: \textbf{Gap~C}.

%----------------------------------------------------------------------
\subsection{The problem: complexification is not algebraically forced}
\label{sec:intro-problem}
%----------------------------------------------------------------------

Paper~7's gap analysis~\cite{Paper7} showed that four independent
algebraic routes to complexification all fail.  The
obstruction is structural: the Peirce eigenspace
$V_1 = \mathbb{R} \cdot E_{11}$ is one-dimensional, creating a
bottleneck through which no Peirce-mediated mechanism can transmit a
complex structure from the observer's C*-algebra to the half-space
$\Vhalf$.  This result is sharp.  Complexification of $\Vhalf$ is
\emph{not} forced by the Jordan algebra structure of $\hthree$, by the
Peirce decomposition, or by any known algebraic property of the
observer--universe interface~\cite{BGW2020}.

The failure is not a computational gap awaiting a cleverer argument.
It is a no-go result: the $V_1$ bottleneck is a theorem about the
Peirce decomposition, and it closes the algebraic route to Gap~C.
Any resolution must come from outside the algebra.

%----------------------------------------------------------------------
\subsection{The resolution: thermodynamic selection}
\label{sec:intro-resolution}
%----------------------------------------------------------------------

We narrow Gap~C via two independent chains of results that
converge on a common geometric prerequisite.

\medskip\noindent
\textbf{Chain~1 (all self-modelers):}
The Landauer bound~\cite{Landauer1961,Bennett1982} applied to the
self-modeling update cycle (Sec.~\ref{sec:landauer}) shows that
maintaining correlations $\MI{B}{M} > 0$ dissipates at least
$\kB T\,\MI{B}{M}$ per cycle.  Free energy requires
non-equilibrium ($F > F_{\mathrm{eq}}$), and a finite system out of
equilibrium must have $S(t) < \Smax$.  The entropy gradient theorem
(Sec.~\ref{sec:egt}), proved via three independent routes, establishes
this for \emph{all} self-modelers regardless of their matter content.
The entropy gradient defines a temporal direction---the arrow of time.

\medskip\noindent
\textbf{Chain~2 (SM-like observers):}
Chirality requires complexification: the $\Cl{6}$ volume form
$\omegasix$ has no eigenstates in real
$\Vhalf = \mathbb{O}^2$~\cite{Paper7}.  Chirality also requires
time-orientation: the chirality operator $\Gammastar$ in Lorentzian
signature flips under time reversal~\cite{LawsonMichelsohn1989}
(Sec.~\ref{sec:chirality}).  Therefore SM-like observers require
both complexification and time-orientation.

\medskip\noindent
\textbf{Structural convergence:}
Both chains require time-orientation---Chain~1 for the thermodynamic
arrow, Chain~2 for the Weyl decomposition.  The single algebraic
choice $u \in S^6$ that determines gauge group and
chirality~\cite{Paper7} simultaneously creates a requirement for the
same time-orientation that the entropy gradient provides.  This
convergence is a structural consistency check---the framework's
particle physics and thermodynamics both require a geometric
property that the emergent spacetime provides.  It is \emph{not}
a logical closure:
the two chains are independent, and we do not prove that
non-complexified blocks have $\rhoexp = 0$.

%----------------------------------------------------------------------
\subsection{What this paper proves and does not prove}
\label{sec:intro-contract}
%----------------------------------------------------------------------

\noindent\textbf{Proved (within stated assumptions):}
\begin{itemize}
  \item \textbf{Entropy gradient theorem}
    (Theorem~\ref{thm:entropy-gradient}): every self-modeling composite
    with $\rhoexp > 0$ satisfies $S(t) < \Smax$, established via three
    independent routes using different assumption subsets
    (Sec.~\ref{sec:egt}).
  \item \textbf{Landauer bound on self-modeling}
    (Theorem~\ref{thm:landauer}): each update cycle dissipates at
    least $\kB T\,\MI{B}{M}$ of work (Sec.~\ref{sec:landauer-bound}).
  \item \textbf{Three-consequence theorem}
    (Theorem~\ref{thm:three-consequence}): the single choice
    $u \in S^6$ determines the gauge group, chirality, \emph{and}
    the requirement for spacetime time-orientation---extending
    Paper~7's ``one choice, two consequences'' to three
    (Sec.~\ref{sec:three-consequences}).
  \item \textbf{Gap~C narrowing}
    (Theorem~\ref{thm:complexification}): complexification is necessary
    for SM-like observers; the status of non-SM self-modelers in
    non-complexified blocks remains open
    (Sec.~\ref{sec:gapc-theorem}).
\end{itemize}

\medskip
\noindent\textbf{Not proved:}
\begin{itemize}
  \item \textbf{Gap~C is narrowed, not closed.}  The $V_1$
    bottleneck remains genuine~\cite{BGW2020}.  Non-complexified blocks
    exist and are not forbidden by the algebra.  Whether they can
    support non-SM self-modelers is an open question.
  \item \textbf{The Past Hypothesis is not derived.}  The entropy
    gradient theorem shows that self-modeling \emph{requires}
    $S(t) < \Smax$ but does not explain \emph{why} the initial entropy
    was low~\cite{Penrose1979,Albert2000}.  Its status is elevated from
    ``observed cosmological fact'' to ``necessary condition for
    self-modeling''---an observer-selection argument, not a
    first-principles derivation.
  \item \textbf{The Second Law is not derived.}  It is \emph{used} as
    input (standard physics) in Link~2 of the chain.
  \item \textbf{The spectral action is not computed.}  The dynamics
    (GR~+~SM Lagrangian from the Connes--Chamseddine--Marcolli spectral
    triple~\cite{CCM2007}) remains a separate open problem.
\end{itemize}

\medskip
\noindent\textbf{Epistemic honesty.}
Gap~C is \emph{narrowed}, not \emph{closed}.  ``Narrowed'' means
we prove complexification is necessary for SM-like observers, but
do not prove $\rhoexp = 0$ for all non-complexified blocks.  A
non-complexified block could in principle have time-orientation
(from dynamics), an entropy gradient, and non-SM self-modelers with
vector-like rather than chiral matter.  Whether such physics supports
self-modeling is open.  This does not undermine the program: it
explains \emph{our} physics (we observe SM chirality, which requires
complexification), not all possible physics---the same epistemic
status as ``why 3+1 dimensions?''

%----------------------------------------------------------------------
\subsection{Paper organization}\label{sec:intro-outline}
%----------------------------------------------------------------------

Section~\ref{sec:entropy} derives the entropy dynamics of self-modeling
composites under the SWAP Hamiltonian: the CPTP channel, its unitality,
and three regimes of entropy behavior (oscillation, monotonic increase,
effective thermalization).

Section~\ref{sec:landauer} establishes the Landauer bound on
self-modeling.  The self-modeling update cycle maps onto Bennett's
Maxwell's demon analysis~\cite{Bennett1982}, and the Reeb--Wolf quantum
Landauer bound~\cite{ReebWolf2014} gives the minimum dissipation per
cycle.  Three independent arguments close a potential coherence
loophole.

Section~\ref{sec:chirality} shows that physical realization of the
chirality from Paper~7 requires the emergent spacetime to carry a
time-orientation, extending ``one choice, two consequences'' to
``one choice, three consequences.''

Section~\ref{sec:gradient-gapc} brings the preceding results together.
The entropy gradient theorem is proved via three independent routes.
Two independent chains---one thermodynamic, one algebraic---converge
on time-orientation, narrowing Gap~C\@.  The master theorem collects
the results, and an assumption register classifies every input by
severity.

Section~\ref{sec:predictions} extracts quantitative predictions:
the minimum entropy deficit for self-modeling, comparison with
Penrose's gravitational entropy estimate~\cite{Penrose1979} (a
94-order-of-magnitude gap that is a feature, not a failure), the
exhaustion timescale, and the experiential density profile.

Section~\ref{sec:discussion} places the results in the context of
the self-modeling program (Papers~5--7) and discusses connections to
the spectral triple approach~\cite{CCM2007}, Jacobson's thermodynamic
gravity~\cite{Jacobson1995}, and observer-selection
principles~\cite{Hoffman2015}.

Appendix~\ref{app:derivations} collects detailed derivations, and
Appendix~\ref{app:assumptions} provides the complete assumption
register.
